\chapter{Basis of the metric Vandermonde matrix}
\label{ch:metric-basis}

% Line colours of the plots in this chapter (monomial, Legendre, Chebyshev,
% Lagrange, ...)
\pgfplotscreateplotcyclelist{metricbasis}{%
  red\\blue\\black\\green!50!black\\orange\\violet\\}

The metric (trial) shape functions are computed from a generalized Vandermonde
matrix, i.e.\ from the values of a set of basis polynomials at the element nodes
(\cref{ch:theory}). Any basis of the polynomial space can be used for this
purpose. This chapter summarizes the implementation of the four bases compared in
the paper and of the hybrid trial space (\code{method=2}). The results are those of
\code{manuscript/studies/ms_element_study.m} (data:
\code{manuscript/data/element-*.dat}); the discussion is given in the paper. The
tests are in \code{tests/unit/MetricBasisTest.m} and
\code{tests/unit/MetricHybridTest.m}.

\section{Setting and basis invariance}

An element of degree $(p_\upxi,p_\upeta)$ has $n=(p_\upxi+1)(p_\upeta+1)$ nodes
$\vect x_k$ with skew coordinates $(s_k,t_k)$. Its metric shape functions $M_l$,
$l=1,\dots,n$, are polynomials of the space $\mathcal Q_p(s,t)$ (degree $p_\upxi$ in $s$
and $p_\upeta$ in $t$, \cref{ch:theory}) and are interpolatory,
$M_l(\vect x_k)=\delta_{lk}$. To compute them, the code takes $n$ basis
polynomials $\phi_c$ of $\mathcal Q_p(s,t)$, evaluates them at the nodes,
$V_{ck}=\phi_c(\vect x_k)$, and at the query points (Gauss points),
$P_{cq}=\phi_c(\vect x_q)$, and solves
\[
  \matr V\,\matr N=\matr P,\qquad N_{lq}=M_l(\vect x_q).
\]

\paragraph{The basis does not change the shape functions.} If $\matr V$ is
non-singular, $\mathcal Q_p(s,t)$ contains exactly one set of functions with the
Kronecker property $M_l(\vect x_k)=\delta_{lk}$, so every basis leads to the same
$M_l$. (A change of basis multiplies $\matr V$ and $\matr P$ from the left with
the same non-singular matrix, which cancels in $\matr V^{-1}\matr P$.)

\paragraph{The basis changes the round-off error.} The system is solved by an LU
decomposition with partial pivoting. The relative error of the computed shape
functions is then bounded, up to a modest factor that depends on $n$, by
$\kappa(\matr V)\,\epsilon$, where $\epsilon\approx1.1\times10^{-16}$ is the unit
round-off of double precision \parencite{Higham2002Accuracy}. Roughly speaking,
every decade of the condition number $\kappa(\matr V)$ costs one significant
digit.

\section{Bases}

\code{PolynomialBasis1D(s, n, basis)} provides, for the polynomial
degree $p$ (argument \code{n}), with values and first derivatives,
\begin{center}
  \begin{tabularx}{\linewidth}{@{}lL@{}}
    \toprule
    \code{basis} & $\phi_a(s)$, $a=0,\dots,p$ \\
    \midrule
    \code{'monomial'} & $s^a$ (used without scaling) \\
    \code{'legendre'} & Legendre polynomials $L_a(s)$, three-term recurrence \\
    \code{'chebyshev'} & Chebyshev polynomials $T_a(s)$, three-term recurrence \\
    \code{'lagrange'} & Lagrange polynomials $\ell_a(s)$ through the 1D reference
      nodes \code{nodal_distribution(n)}, product formula (no division by
      the differences $s-s_m$ to the nodes $s_m$, so query points may coincide
      with nodes) \\
    \bottomrule
  \end{tabularx}
\end{center}
The 2D basis is the tensor product $\phi_a(\tilde s)\,\phi_b(\tilde t)$. The
Legendre, Chebyshev and Lagrange bases are defined on $[-1,1]$, but on distorted
elements the skew coordinates of the nodes leave $[-1,1]^2$. Outside $[-1,1]$ these
polynomials grow quickly. For these three bases \code{ShapeFunctions2D_Metric}
therefore maps the bounding box of the nodes in $(s,t)$ onto $[-1,1]^2$,
\begin{equation}
  \tilde s = (s-c_\mathrm{s})/r_\mathrm{s},\qquad
  \tilde t = (t-c_\mathrm{t})/r_\mathrm{t},
\end{equation}
where $(c_\mathrm{s},c_\mathrm{t})$ is the centre and $(r_\mathrm{s},r_\mathrm{t})$
are the half widths of the box. The map is affine and computed from
frame-invariant quantities, so the space $\mathcal Q_p$ and the frame invariance are
preserved; on a parallelogram it is the identity. The monomial basis is used
without this scaling. The basis is selected with \code{p.metricBasis} (default
\code{'lagrange'}).

\Cref{fig:metric-basis-kappa,fig:metric-basis-leb} show the conditioning and the
Lebesgue constants. On a parallelogram, the Lagrange basis gives $\kappa_2(\matr V)=1$
for every $p$, since $\matr V$ is then a permutation matrix. On distorted elements,
the three scaled bases are close to one another, and the Lebesgue constant of
$\mathcal Q_p(s,t)$ grows exponentially with $p$. Since the shape functions do not
depend on the basis, no basis can remove this growth.

\begin{figure}[htbp]
  \centering
  \begin{tikzpicture}
    \begin{groupplot}[group style={group size=2 by 2, horizontal sep=2.6cm,
        vertical sep=2.2cm}, width=0.44\linewidth, height=0.35\linewidth,
        ymode=log, grid=major, xlabel={$p$}, ylabel={$\kappa_2(\matr V)$},
        ymin=0.5, ymax=1e16, xmin=1, xmax=16, cycle list name=metricbasis,
        unbounded coords=discard,
        every axis plot/.append style={mark size=1.4pt}]
      \nextgroupplot[title={parallelogram},
        legend to name=leg:metricbasis, legend columns=4,
        legend style={font=\footnotesize}]
      \addplot+[mark=*] table[x=p, y=kappa2_mon]{../manuscript/data/element-parallelogram.dat};
      \addplot+[mark=square*] table[x=p, y=kappa2_leg]{../manuscript/data/element-parallelogram.dat};
      \addplot+[mark=triangle*] table[x=p, y=kappa2_che]{../manuscript/data/element-parallelogram.dat};
      \addplot+[mark=diamond*] table[x=p, y=kappa2_lag]{../manuscript/data/element-parallelogram.dat};
      \legend{monomial, Legendre, Chebyshev, Lagrange}
      \nextgroupplot[title={quadrilateral}]
      \addplot+[mark=*] table[x=p, y=kappa2_mon]{../manuscript/data/element-trapezoid.dat};
      \addplot+[mark=square*] table[x=p, y=kappa2_leg]{../manuscript/data/element-trapezoid.dat};
      \addplot+[mark=triangle*] table[x=p, y=kappa2_che]{../manuscript/data/element-trapezoid.dat};
      \addplot+[mark=diamond*] table[x=p, y=kappa2_lag]{../manuscript/data/element-trapezoid.dat};
      \nextgroupplot[title={curved}]
      \addplot+[mark=*] table[x=p, y=kappa2_mon]{../manuscript/data/element-curved.dat};
      \addplot+[mark=square*] table[x=p, y=kappa2_leg]{../manuscript/data/element-curved.dat};
      \addplot+[mark=triangle*] table[x=p, y=kappa2_che]{../manuscript/data/element-curved.dat};
      \addplot+[mark=diamond*] table[x=p, y=kappa2_lag]{../manuscript/data/element-curved.dat};
      \nextgroupplot[title={general}]
      \addplot+[mark=*] table[x=p, y=kappa2_mon]{../manuscript/data/element-general.dat};
      \addplot+[mark=square*] table[x=p, y=kappa2_leg]{../manuscript/data/element-general.dat};
      \addplot+[mark=triangle*] table[x=p, y=kappa2_che]{../manuscript/data/element-general.dat};
      \addplot+[mark=diamond*] table[x=p, y=kappa2_lag]{../manuscript/data/element-general.dat};
    \end{groupplot}
  \end{tikzpicture}\\[4pt]
  \ref{leg:metricbasis}
  \caption{Condition number of the metric Vandermonde matrix on the single test
    elements of \cref{fig:test-elements} (bounding-box scaling for the Legendre,
    Chebyshev and Lagrange bases).}
  \label{fig:metric-basis-kappa}
\end{figure}

\section{Hybrid trial space (\texttt{method=2})}
\label{sec:metric-hybrid}

\paragraph{Definition.} With $m=\min(p_\upxi,p_\upeta)$ and the Legendre
polynomials $L_a$, the hybrid trial space is
\begin{equation}
  \mathcal T_p = \mathcal P_m(\tilde s,\tilde t)\;\oplus\;\operatorname{span}\{L_a(\xi)L_b(\eta):\
  a\le p_\upxi,\ b\le p_\upeta,\ a+b>m\}.
  \label{eq:hybrid-space}
\end{equation}
The first part, $\mathcal P_m(\tilde s,\tilde t)$, consists of all polynomials of total
degree $\le m$ in the bounding-box scaled skew coordinates $(\tilde s,\tilde t)$;
the code spans it with the products $L_a(\tilde s)L_b(\tilde t)$, $a+b\le m$.
Because $\tilde s$ and $\tilde t$ are affine functions of $x$ and $y$, and vice
versa, $\mathcal P_m(\tilde s,\tilde t)=\mathcal P_m(x,y)$. The hybrid space therefore
contains all complete polynomials of degree $m$ in $x$ and $y$ and, since
$(\tilde s,\tilde t)$ move and rotate with the element, it is frame invariant. The
second part, the complement, is written in the parametric coordinates $(\xi,\eta)$,
in which the element nodes form a tensor grid. On non-affine elements the complement
depends on the polynomials used. The Legendre basis is therefore part of the
definition, and \code{p.metricBasis} is ignored for \code{method=2}.

\paragraph{Straight-edged elements.} If the geometry map is bilinear and the
hybrid Vandermonde matrix is non-singular, the hybrid space equals the
parametric space, $\mathcal T_p=\mathcal Q_p(\xi,\eta)$: the metric functions coincide with the
parametric ones, and the element is the classical (symmetric) spectral element.
The reason is that $x$ and $y$ are bilinear in $(\xi,\eta)$ for such a map, so a
polynomial of total degree $\le m$ in $x$ and $y$ has degree $\le m$ in $\xi$ and
in $\eta$ \parencite[cf.][]{Arnold2002Approximation}.

\paragraph{Folded and inverted elements.} The hybrid functions are evaluated only
through $(\xi,\eta)$, so the sign of $\det(\matr J)$ is irrelevant, and completeness
holds pointwise in $(\xi,\eta)$, also where the map folds. Only $\det(\matr J)=0$ at
an evaluation point is excluded, because the gradients of the complement contain
$\matr J^{-1}$ (error \code{ShapeFunctions2D_Metric:singularJacobian}).

\paragraph{Implementation.} The hybrid space is selected with
\code{p.metricMethod=2} or for a single call with
\code{ShapeFunctions2D_Metric(..., p, 2, xiQ)}. The parametric coordinates of the
nodes are the reference nodes (\code{ReferenceNodes2D}). Those of the query points
are passed by the callers, which evaluate at mapped quadrature or sampling points.
The gradients of $\mathcal P_m(\tilde s,\tilde t)$ follow from the constant matrix
$\matr J_0^{-1}$ of the skew coordinates, those of the complement from the Jacobian
at the query point, $\nabla_{\vect x}\phi=\matr J^{-1}\nabla_{\xi}\phi$.

\begin{figure}[htbp]
  \centering
  \begin{tikzpicture}
    \begin{semilogyaxis}[width=0.62\linewidth, height=0.42\linewidth,
        xlabel={$p$}, ylabel={Lebesgue constant}, grid=major, ymax=1e12,
        legend cell align=left, legend style={font=\footnotesize, at={(0.02,0.98)}, anchor=north west},
        cycle list name=metricbasis]
      \addplot+[mark=square*] table[x=p, y=LQ]{../manuscript/data/element-trapezoid.dat};
      \addplot+[mark=triangle*] table[x=p, y=LQ]{../manuscript/data/element-general.dat};
      \addplot+[mark=o] table[x=p, y=LQ]{../manuscript/data/element-curved.dat};
      \addplot+[mark=square, dashed] table[x=p, y=LT]{../manuscript/data/element-trapezoid.dat};
      \addplot+[mark=triangle, dashed] table[x=p, y=LT]{../manuscript/data/element-general.dat};
      \addplot+[mark=x, dashed] table[x=p, y=LT]{../manuscript/data/element-curved.dat};
      \legend{{$\mathcal Q_p(s,t)$ quadrilateral}, {$\mathcal Q_p(s,t)$ general}, {$\mathcal Q_p(s,t)$ curved},
              $\mathcal T_p$ quadrilateral, $\mathcal T_p$ general, $\mathcal T_p$ curved}
    \end{semilogyaxis}
  \end{tikzpicture}
  \caption{Lebesgue constant of the metric interpolation in $\mathcal Q_p(s,t)$
    (\texttt{method=1}) and in the hybrid space $\mathcal T_p$ (\texttt{method=2}) on
    the distorted test elements of \cref{fig:test-elements}.}
  \label{fig:metric-basis-leb}
\end{figure}
